\subsection{Symmetric polynomials}

Let \(n\) be a positive integer. For every permutation \(\sigma \in \mathfrak{S}_n\) let \(\varphi_{\sigma}\) be the automorphism of the A-algebra \(\mathbf{A}[\mathbf{X}_1, \ldots, \mathbf{X}_n]\) which maps \(\mathbf{X}_i\) to \(\mathbf{X}_{\sigma(i)}\) for \(1 \leqslant i \leqslant n\) . It is clear that \(\sigma \mapsto \varphi_{\sigma}\) is a homomorphism of \(\mathfrak{S}_n\) into the group of automorphisms of \(\mathbf{A}[\mathbf{X}_1, \ldots, \mathbf{X}_n]\) . We shall put \(\sigma f = \varphi_{\sigma}(f)\) for \(\sigma \in \mathfrak{S}_n\) and \(f \in \mathbf{A}[\mathbf{X}_1, \ldots, \mathbf{X}_n]\) . The polynomial \(f\) is said to be symmetric if \(\sigma f = f\) for all \(\sigma \in \mathfrak{S}_n\) ; the symmetric polynomials form a unital graded subalgebra of \(\mathbf{A}[\mathbf{X}_1, \ldots, \mathbf{X}_n]\) ; we shall denote it by \(\mathbf{A}[\mathbf{X}_1, \ldots, \mathbf{X}_n]^{\text{sym}}\) in the rest of this paragraph.

For every positive integer \(k\) we denote by \(\mathfrak{P}_k\) the set of \(k\) -element subsets of the set \(\{1, 2, \ldots, n\}\) and put

\[
s _ {k} = \sum_ {\mathrm{H} \in \mathfrak {P} _ {k}} \prod_ {i \in \mathrm{H}} X _ {i}.\tag{1}
\]

When we wish to specify the integer \(n\) we shall write \(s_{k,n}\) instead of \(s_k\) . We have in particular

\[
\begin{array}{l} s _ {0} = 1 \\ s _ {1} = \sum_ {1 \leqslant i \leqslant n} X _ {i} \\ s _ {2} = \sum_ {1 \leqslant i <   j \leqslant n} X _ {i} X _ {j} \\ \dots \dots \\ s _ {n} = X _ {1} \dots X _ {n} \end{array}
\]

and \(s_k = 0\) for \(k > n\) . It is clear that \(s_k\) is a homogeneous symmetric polynomial of degree \(k\) ; we shall call it the elementary symmetric polynomial of degree \(k\) .

In the ring \(\mathbf{A}[\mathbf{X}_1,\dots ,\mathbf{X}_n,\mathbf{U},\mathbf{V}]\) we have the relation

\[
\prod_ {i = 1} ^ {n} \left(\mathrm{U} + \mathrm{VX} _ {i}\right) = \sum_ {k = 0} ^ {n} \mathrm{U} ^ {n - k} \mathrm{V} ^ {k} s _ {k};\tag{2}
\]

by appropriate substitutions we deduce the relations

\[
\prod_ {i = 1} ^ {n} \left(1 + \mathrm{TX} _ {i}\right) = \sum_ {k = 0} ^ {n} s _ {k} \mathrm{T} ^ {k},\tag{3}
\]

\[
\prod_ {i = 1} ^ {n} \left(\mathbf {X} - \mathbf {X} _ {i}\right) = \sum_ {k = 0} ^ {n} (- 1) ^ {n - k} s _ {n - k} \mathbf {X} ^ {k}.\tag{4}
\]

THEOREM 1. --- Let us put \(E = A{[}X_{1}, \ldots, X_{n}{]}\) and \(S = A{[}X_{1}, \ldots, X_{n}{]}^{\text{sym}}\).

\begin{enumerate}
\def\labelenumi{\alph{enumi})}
\item
  The A-algebra S of symmetric polynomials is generated by \(s_1, \ldots, s_n\) .
\item
  The elements \(s_1, \ldots, s_n\) of \(\mathbf{E}\) are algebraically independent over \(\mathbf{A}\) (IV, p.~4).
\item
  The family of monomials \(\mathbf{X}^{\nu} = \mathbf{X}_1^{\nu(1)} \ldots \mathbf{X}_n^{\nu(n)}\) such that \(0 \leqslant \nu(i) < i\) for \(1 \leqslant i \leqslant n\) is a basis of the S-module \(\mathbf{E}\) . In particular, \(\mathbf{E}\) is a free S-module of rank \(n!\) .
\end{enumerate}

We shall prove the theorem by induction on \(n\) , the case \(n = 0\) being trivial. Let us write \(\mathbf{B} = \mathbf{A}[\mathbf{X}_n]\) and denote by \(s_k'\) the elementary symmetric polynomial of degree \(k\) in \(\mathbf{X}_1, \ldots, \mathbf{X}_{n-1}\) ; we thus have \(\mathbf{B}[\mathbf{X}_1, \ldots, \mathbf{X}_{n-1}] = \mathbf{A}[\mathbf{X}_1, \ldots, \mathbf{X}_n]\) . If we replace \(n\) by \(n-1\) and \(\mathbf{A}\) by \(\mathbf{B}\) in the statement of Th. 1, we can formulate the induction hypothesis as follows:

\begin{enumerate}
\def\labelenumi{(\Alph{enumi})}
\item
  The B-algebra \(S'\) of polynomials \(f \in A[X_1, \ldots, X_n]\) invariant under all permutations of \(X_1, \ldots, X_{n-1}\) is generated by \(s_1', \ldots, s_{n-1}'\) .
\item
  The elements \(s_1', \ldots, s_{n-1}'\) of \(E\) are algebraically independent over \(B\) .
\item
  The family of monomials \(\mathbf{X}_{1}^{\nu(1)}\ldots\mathbf{X}_{n-1}^{\nu(n-1)}\) such that \(0\leqslant\nu(i)<i\) for \(1\leqslant i\leqslant n-1\) is a basis of the S'-module E.
\end{enumerate}

We have the obvious relation

\[
s _ {k} = s _ {k} ^ {\prime} + s _ {k - 1} ^ {\prime} \mathbf {X} _ {n} \quad (1 \leqslant k \leqslant n - 1),\tag{5}
\]

from which we obtain by induction on \(k\)

\[
s _ {k} ^ {\prime} = (- 1) ^ {k} \mathbf {X} _ {n} ^ {k} + \sum_ {i = 1} ^ {k} (- 1) ^ {k - i} s _ {i} \mathbf {X} _ {n} ^ {k - i} \quad (1 \leqslant k \leqslant n - 1).\tag{6}
\]

We have \(\mathbf{S} \subset \mathbf{S}'\) , hence \(s_1, \ldots, s_n\) belong to \(\mathbf{S}'\) ; by (A) and formula (6) the B-algebra \(\mathbf{S}'\) is therefore generated by \(s_1, \ldots, s_{n-1}\) .

By (B) there exists an endomorphism \(u\) of the B-algebra \(S'\) such that

\[
u \left(s _ {k} ^ {\prime}\right) = (- 1) ^ {k} X _ {n} ^ {k} + \sum_ {i = 1} ^ {k} (- 1) ^ {k - i} s _ {i} ^ {\prime} X _ {n} ^ {k - i} \quad (1 \leqslant k \leqslant n - 1).\tag{7}
\]

By (5), we have \(u(s_k) = u(s_k') + u(s_{k-1}') \mathbf{X}_n\) , whence \(u(s_k) = s_k'\) by an easy calculation. Let \(\mathbf{P} \in \mathbf{B}[\mathbf{Y}_1, \ldots, \mathbf{Y}_{n-1}]\) ; then from \(\mathbf{P}(s_1, \ldots, s_{n-1}) = 0\) we deduce

\[
0 = u \left(\mathrm{P} \left(s _ {1}, \dots , s _ {n - 1}\right)\right) = \mathrm{P} \left(s _ {1} ^ {\prime}, \dots , s _ {n - 1} ^ {\prime}\right),
\]

whence P = 0 by (B). It follows that the B-algebra \(S'\) is generated by the algebraically independent elements \(s_{1}, \ldots, s_{n-1}\) . This property may be reformulated as follows:

\begin{enumerate}
\def\labelenumi{(\Alph{enumi})}
\setcounter{enumi}{3}
\tightlist
\item
  The A-algebra \(S'\) is generated by the algebraically independent elements \(s_1, \ldots, s_{n-1}, X_n\) .
\end{enumerate}

We can thus identify \(S'\) with the polynomial ring \(C[X_n]\) , where \(C\) is the sub-A-algebra of \(E\) generated by \(s_1, \ldots, s_{n-1}\) .

To prove a), let \(f \in S\) be a homogeneous symmetric polynomial of degree \(m\) . We have \(f \in S' = C[X_n]\) , so there exists an element \(g = P(s_1, \ldots, s_{n-1})\) of \(C\) , homogeneous of degree \(m\) in \(X_1, \ldots, X_n\) , such that \(f - g\) is divisible by \(X_n\) . Since \(f - g\) is symmetric, each of its terms is also divisible by \(X_1, \ldots, X_{n-1}\) , hence \(f - g\) is divisible by \(s_n = X_1 \ldots X_n\) . In other words, there exists \(h \in S\) such that \(f = g + hs_n\) , whence \(\deg h < m\) . By induction on \(m\) it follows that \(f\) belongs to

\[
\mathbf {C} \left[ s _ {n} \right] _ {\mathrm{E}} = \mathbf {A} \left[ s _ {1}, \dots , s _ {n - 1}, s _ {n} \right] _ {\mathrm{E}}.
\]

Thus we have shown that the A-algebra S is generated by \(s_{1}, \ldots, s_{n}\) .

Next we prove \(b\) ). If we substitute \(\mathbf{X}_n\) for \(\mathbf{X}\) in (4), we find

\[
(- 1) ^ {n + 1} s _ {n} = \mathrm{X} _ {n} ^ {n} + \sum_ {k = 1} ^ {n - 1} (- 1) ^ {n - k} s _ {n - k} \mathrm{X} _ {n} ^ {k};
\]

in other words, \((-1)^{n + 1}s_n\) is a monic polynomial in \(X_{n}\) of degree \(n\) and with coefficients in C. By IV, p.~11 we thus have the following property:

\begin{enumerate}
\def\labelenumi{(\Alph{enumi})}
\setcounter{enumi}{4}
\tightlist
\item
  The C-algebra homomorphism \(\varphi\) of C{[}T{]} into C{[}X \(_n\) {]} = S' which maps T to \(s_n\) is injective, and S' is a free module over the image of \(\varphi\) , with basis (1, X \(_n\) , \ldots, X \(_n^{n-1}\) ).
\end{enumerate}

The elements \(s_{1}, \ldots, s_{n-1}\) of C are algebraically independent over A, by (D); thus the injectivity of \(\varphi\) means that \(s_{1}, \ldots, s_{n-1}, s_{n}\) are algebraically independent over A, whence b).

To prove c), the image of \(\varphi\) equals \(C[s_{n}]_{E} = S\) , hence by (E), \(S'\) is a free module over S on the basis \((1, X_{n}, ..., X_{n}^{n-1})\) . Now the assertion c) follows from the induction hypothesis (C) and Prop. 25 of II, p.~222. This completes the proof.

Let \(f\) be a symmetric polynomial in \(\mathbf{X}_1, \ldots, \mathbf{X}_n\) , homogeneous of degree \(m\) . By Th. 1 (IV, p.~62) there exists a polynomial \(\mathbf{Q} \in \mathbf{A}[\mathbf{Y}_1, \ldots, \mathbf{Y}_n]\) such that \(f = \mathbf{Q}(s_1, \ldots, s_n)\) . The preceding proof provides a means of explicit calculation of \(\mathbf{Q}\) , by a double induction on \(n\) and \(m\) . For as we have seen, there exists a polynomial \(\mathbf{P} \in \mathbf{A}[\mathbf{Y}_1, \ldots, \mathbf{Y}_{n-1}]\) and a polynomial \(h\) symmetric in \(\mathbf{X}_1, \ldots, \mathbf{X}_n\) which is homogeneous of degree \(m - n\) , such that

\[
f = \mathrm{P} (s _ {1}, \dots , s _ {n - 1}) + s _ {n} h.\tag{8}
\]

For every polynomial \(u \in \mathbf{A}[\mathbf{X}_1, \ldots, \mathbf{X}_n]\) we put

\[
u ^ {\prime} \left(\mathrm{X} _ {1}, \dots , \mathrm{X} _ {n - 1}\right) = u \left(\mathrm{X} _ {1}, \dots , \mathrm{X} _ {n - 1}, 0\right).
\]

Then \(s_1', \ldots, s_{n-1}'\) are the elementary symmetric polynomials in \(X_1, \ldots, X_{n-1}\) and formula (8) implies

\[
f ^ {\prime} = \mathrm{P} \left(s _ {1} ^ {\prime}, \dots , s _ {n - 1} ^ {\prime}\right).
\]

The determination of P is thus reduced to a calculation of symmetric polynomials in n - 1 indeterminates, and h is obtained from (8).

We illustrate the method by two examples.

Examples. --- 1) Let \(n = 3\) and

\[
f = \mathbf {X} _ {1} ^ {2} (\mathbf {X} _ {2} + \mathbf {X} _ {3}) + \mathbf {X} _ {2} ^ {2} (\mathbf {X} _ {3} + \mathbf {X} _ {1}) + \mathbf {X} _ {3} ^ {2} (\mathbf {X} _ {1} + \mathbf {X} _ {2}).
\]

We have

\[
f ^ {\prime} = \mathbf {X} _ {1} ^ {2} \mathbf {X} _ {2} + \mathbf {X} _ {1} \mathbf {X} _ {2} ^ {2} = \mathbf {X} _ {1} \mathbf {X} _ {2} (\mathbf {X} _ {1} + \mathbf {X} _ {2}) = s _ {1} ^ {\prime} s _ {2} ^ {\prime}.
\]

Let us write \(g = f - s_1s_2\) ; we have

\[
\boldsymbol {g} = f - \left(\mathbf {X} _ {1} + \mathbf {X} _ {2} + \mathbf {X} _ {3}\right) \left(\mathbf {X} _ {1} \mathbf {X} _ {2} + \mathbf {X} _ {1} \mathbf {X} _ {3} + \mathbf {X} _ {2} \mathbf {X} _ {3}\right) = - 3 \mathbf {X} _ {1} \mathbf {X} _ {2} \mathbf {X} _ {3},
\]

hence finally

\[
f = s _ {1} s _ {2} - 3 s _ {3}.
\]

\begin{enumerate}
\def\labelenumi{\arabic{enumi})}
\setcounter{enumi}{1}
\tightlist
\item
  Again let \(n = 3\) and put \(p = \mathbf{X}_1^3 + \mathbf{X}_2^3 + \mathbf{X}_3^3\) .
\end{enumerate}

We have \(p(\mathbf{X}_1, 0, 0) = \mathbf{X}_1^3 = s_1(\mathbf{X}_1, 0, 0)^3\) . Writing \(q = p - s_1^3\) , we obtain

\[
q = - 3 f - 6 \mathrm{X} _ {1} \mathrm{X} _ {2} \mathrm{X} _ {3} = - 3 s _ {1} s _ {2} + 3 s _ {3}
\]

and finally

\[
p = s _ {1} ^ {3} - 3 s _ {1} s _ {2} + 3 s _ {3}.
\]

Let \(S_1, \ldots, S_n\) be indeterminates. We shall equip the polynomial algebra \(A[S_1, \ldots, S_n]\) with the graduation of type N for which \(S_k\) is of weight k for \(1 \leqslant k \leqslant n\) (IV, p.~3), and we equip \(A[X_1, \ldots, X_n]\) with the usual graduation. For \(1 \leqslant k \leqslant n\) the elementary symmetric polynomial \(s_{k,n}\) in \(X_1, \ldots, X_n\) is homogeneous of degree k. By Th. 1 (IV, p.~62) the mapping \(g \mapsto g(s_{1,n}, \ldots, s_{n,n})\) is thus an isomorphism of graded algebras

\[
\varphi_ {n}: \mathrm{A} [ \mathrm{S} _ {1}, \dots , \mathrm{S} _ {n} ] \rightarrow \mathrm{A} [ \mathrm{X} _ {1}, \dots , \mathrm{X} _ {n} ] ^ {\text { sym }}.
\]

Let \(m\) be an integer such that \(0 \leqslant m \leqslant n\) . For every integer \(k\) such that \(1 \leqslant k \leqslant m\) we have

\[
s _ {k, m} (\mathbf {X} _ {1}, \dots , \mathbf {X} _ {m}) = s _ {k, n} (\mathbf {X} _ {1}, \dots , \mathbf {X} _ {m}, 0, \dots , 0)\tag{9}
\]

by definition (IV, p.~61, formula (1)) of \(s_k\) . Hence the diagram

\[
\begin{array}{c} \mathrm{A} [ \mathrm{S} _ {1}, \dots , \mathrm{S} _ {m} ] \xrightarrow {j} \mathrm{A} [ \mathrm{S} _ {1}, \dots , \mathrm{S} _ {n} ] \\ \varphi_ {m} \Bigg | _ {\downarrow} \qquad \qquad \qquad \qquad \qquad \qquad \qquad \qquad \qquad \qquad \varphi_ {n} \\ \mathrm{A} [ \mathrm{X} _ {1}, \dots , \mathrm{X} _ {m} ] ^ {\text { sym }} \xleftarrow {\rho} \mathrm{A} [ \mathrm{X} _ {1}, \dots , \mathrm{X} _ {n} ] ^ {\text { sym }} \end{array}\tag{10}
\]

(where \(j\) denotes the canonical inclusion and \(\rho\) the homomorphism

\[
g \mapsto g \left(\mathrm{X} _ {1}, \dots , \mathrm{X} _ {m}, 0, \dots , 0\right))
\]

is commutative.

PROPOSITION 1. --- For every pair of positive integers \(k\) , \(n\) , let \(\mathbf{S}_{k}^{(n)}\) be the A-module consisting of all symmetric polynomials in \(\mathbf{X}_{1},\ldots,\mathbf{X}_{n}\) , which are homogeneous of degree \(k\) . If the integer \(m\) satisfies \(0\leqslant k\leqslant m\leqslant n\) , then the mapping \(f\mapsto f(\mathbf{X}_{1},\ldots,\mathbf{X}_{m},0,\ldots,0)\) is an isomorphism of \(\mathbf{S}_{k}^{(n)}\) onto \(\mathbf{S}_{k}^{(m)}\) .

By the commutativity of the diagram (10) it is enough to show that every isobaric polynomial of weight k in \(S_{1}, \ldots, S_{n}\) depends only on \(S_{1}, \ldots, S_{m}\) under the hypotheses \(0 \leqslant k \leqslant m \leqslant n\) . Now the weight of a monomial \(\mathrm{S}_{1}^{\alpha(1)} \ldots \mathrm{S}_{n}^{\alpha(n)}\) is equal to the integer \(\alpha(1) + 2\alpha(2) + \cdots + n\alpha(n)\) ; since the integers \(\alpha(1), \ldots, \alpha(n)\) are positive, the relation

\[
\alpha (1) + 2 \alpha (2) + \dots + n \alpha (n) = k \leqslant n
\]

implies \(\alpha(j) = 0\) for \(k < j \leqslant n\) , whence the assertion.

Example 3. --- By example 2 of IV, p.64 and Prop. 1 above we thus have

\[
\sum_ {i = 1} ^ {n} X _ {i} ^ {3} = s _ {1, n} ^ {3} - 3 s _ {1, n} s _ {2, n} + 3 s _ {3, n}
\]

for every integer \(n \geqslant 3\) . The commutativity of the diagram (10) moreover gives the formulae

\[
\begin{array}{c} \mathbf {X} _ {1} ^ {3} + \mathbf {X} _ {2} ^ {3} = s _ {1, 2} ^ {3} - 3 s _ {1, 2} s _ {2, 2}, \\ \mathbf {X} _ {1} ^ {3} = s _ {1, 1} ^ {3}. \end{array}
\]

Remark. --- Let n and k be two positive integers. We denote by \(\Delta_{k,n}\) the set of elements of length k in \(N^{n}\) and further equip \(\Delta_{k,n}\) with the order relation, written \(\alpha \preccurlyeq \beta\) , induced by the lexicographic order on \(N^{n}\) (Set Theory, III, p.~157), and define an action of the group \(S_{n}\) on \(N^{n}\) by \((\sigma\alpha)(i) = \alpha(\sigma^{-1}(i))\) for \(\sigma \in S_{n}\) , \(\alpha \in N^{n}\) and \(1 \leqslant i \leqslant n\) . Moreover, we denote by \(D_{k}\) the set of elements \(\alpha = (\alpha(1), \ldots, \alpha(k))\) of \(N^{k}\) such that

\[
\alpha (1) \geqslant \alpha (2) \geqslant \dots \geqslant \alpha (k), \quad \alpha (1) + \dots + \alpha (k) = k.
\]

Suppose that \(k \leqslant n\) and let us identify \(\mathbf{N}^k\) with a subset of \(\mathbf{N}^n\) by the mapping \((\alpha(1), \ldots, \alpha(k)) \mapsto (\alpha(1), \ldots, \alpha(k), 0, \ldots, 0)\) . Then \(D_k\) consists of the elements \(\alpha\) of \(\Delta_{k,n}\) such that \(\sigma \alpha \preccurlyeq \alpha\) for all \(\sigma \in \mathfrak{S}_n\) . It follows that every orbit of the group \(\mathfrak{S}_n\) in \(\Delta_{k,n}\) contains a unique element of \(D_k\) . For every \(\alpha \in D_k\) let \(O(\alpha)\) be the orbit of \(\alpha\) in \(\Delta_{k,n}\) under the operation of \(\mathfrak{S}_n\) ; put

\[
\mathbf {M} (\alpha) = \sum_ {\beta \in O (\alpha)} \mathbf {X} ^ {\beta}.\tag{11}
\]

It follows from Lemma 1 of IV, p.~47, that the family \((\mathbf{M}(\alpha))_{\alpha \in \mathrm{D}_k}\) is a basis of the A-module \(\mathbf{S}_k^{(n)}\) .

For each \(\alpha \in D_k\) put

\[
\mathrm{S} (\alpha) = \prod_ {i = 1} ^ {k} s _ {i} ^ {\alpha (i) - \alpha (i + 1)} \quad (\text { where   it   is   understood   that } \alpha (k + 1) = 0);\tag{12}
\]

since we have \(\sum_{i=1}^{k} i \cdot (\alpha(i) - \alpha(i+1)) = \sum_{i=1}^{k} \alpha(i) = k\) , the symmetric polynomial \(S(\alpha)\) is homogeneous of degree \(k\) . It follows directly from Th. 1 (IV, p.~62) that the family \((S(\alpha))_{\alpha \in D_k}\) is a basis of the A-module \(S_k^{(n)}\) .

Let \(\alpha, \beta\) be in \(D_k\) , and let \(c_{\alpha\beta}\) be the coefficient of the monomial \(X^\beta\) in the polynomial \(S(\alpha)\) defined by (12) in the case \(A = Z\) and \(k = n\) . This is then a positive integer, independent of the ring \(A\) and of the integer \(n\) . By formula (9) (IV, p.~64) we then have

\[
\mathrm{S} (\alpha) = \sum_ {\beta \in \mathrm{D} _ {k}} c _ {\alpha \beta} \cdot \mathrm{M} (\beta) \quad (\alpha \in \mathrm{D} _ {k}).\tag{13}
\]

It may be shown (cf.~IV, p.~101, Exercise 13) that the matrix \(\mathbf{C} = (c_{\alpha \beta})_{\alpha, \beta \in D_k}\) is such that \(c_{\alpha \alpha} = 1\) and \(c_{\alpha \beta} = 0\) for \(\alpha < \beta\) ; generalizing the terminology introduced in II, p.~351, we shall say that \(\mathbf{C}\) belongs to the lower total triangular group. The same is true of the matrix \(\mathbf{D}\) inverse to \(\mathbf{C}\) . The values of the matrices \(\mathbf{C}\) and \(\mathbf{D}\) when \(2 \leqslant k \leqslant 5\) , are given in the table (IV, p.~103-105).

Suppose now that \(n < k\) and let us identify \(\mathbf{N}^n\) with a subset of \(\mathbf{N}^k\) by the mapping \((\alpha(1), \ldots, \alpha(n)) \mapsto (\alpha(1), \ldots, \alpha(n), 0, \ldots, 0)\) . For every \(\alpha\) in \(D_k \cap N^n\) we again denote by \(O(\alpha)\) the orbit of \(\alpha\) in \(\Delta_{k,n}\) under the operation of \(\mathfrak{S}_n\) and define \(M(\alpha)\) by (11). Further we define \(S(\alpha)\) by (12), then the families \((M(\alpha))_{\alpha \in D_k \cap N^n}\) and \((S(\alpha))_{\alpha \in D_k \cap N^n}\) are bases of the A-module \(S_k^{(n)}\) . By formula (9) of IV, p.~64, we have a formula analogous to (13) where \(D_k\) is replaced by \(D_k \cap N^n\) , with the same integers \(c_{\alpha\beta}\) .

Example 4. --- By Example 3 of IV, p.~65, we have

\[
\mathbf {M} (3, 0, 0) = \mathbf {S} (3, 0, 0) - 3 \mathbf {S} (2, 1, 0) + 3 \mathbf {S} (1, 1, 1)
\]

for every integer \(n \geqslant 3\) , and so

\[
\mathbf {M} (3, 0) = \mathbf {S} (3, 0) - 3 \mathbf {S} (2, 1)
\]

for n = 2.

